\section{Problem Definition}
\label{sec:problem}

The objective is to produce, automatically and reproducibly, a weekly timetable for the Department
of Computer Engineering that satisfies every mandatory institutional rule and violates as few
preferential rules as possible. A comprehensive review of the literature confirms that constraint
programming is the predominant methodology for this class of problem, owing to its ability to model
mandatory and preferential requirements in a single framework~\cite{schaerf1999survey,muller2007}.

\subsection{Scheduling model}

In the CPSolver course-timetabling model the decision variable is a \emph{class} (a single lecture,
laboratory or tutorial section), and its value is a \emph{placement}: a starting time together with
the room or rooms in which the class meets. Students are not decision variables. A student's
timetable is a consequence of the placements chosen for the classes that student is enrolled in, and
the cost of two classes overlapping for a shared cohort is priced through a joint-enrolment
constraint between those two classes. This distinction matters for everything that follows: the
solver can only improve a student's timetable by moving a class, or---if it is permitted to do
so---by moving the student into a different section. Student sectioning is otherwise a problem in its own right~\cite{muller2010sectioning}.

\subsection{Hard and soft constraints}

Hard constraints must hold in any acceptable timetable:

\begin{itemize}
  \item no room hosts two classes at overlapping times;
  \item no instructor teaches two classes at overlapping times;
  \item a class is placed only in a room whose capacity meets its enrolment limit and whose type
    and equipment suit the subpart;
  \item every class respects the meeting pattern declared for its subpart (for example
    $2\times120$ minutes) and the date pattern of the term;
  \item rooms are used only when available, and weekends are excluded.
\end{itemize}

Soft constraints express quality rather than legality: preferred rooms and time windows, balanced
instructor load, limited back-to-back teaching, minimised travel between buildings, a protected
lunch window, and---the criterion this report treats as primary---the number of students whose
personal timetable contains an overlap. Soft requirements enter as weighted criteria rather than as further hard constraints~\cite{rudova2003soft}.

\subsection{Student conflicts and the measurement problem}

Two classes conflict for a student if and only if they share a day, their slot intervals overlap,
and, when both carry date patterns, they share at least one active week. Intervals that merely touch
at an endpoint do not conflict, and room and instructor are irrelevant to the test. This rule is
implemented identically in three places that were used during the project---CPSolver's own time
representation, the Java conflict oracle written for the learning layer, and the Python evaluator
described in Section~\ref{sec:evaluation}---and the agreement between them is reported in
Section~\ref{sec:evaluation}.

The central difficulty of this phase was not producing a feasible timetable; coverage, room and
instructor feasibility were reached early. The difficulty was that student-conflict counts could not
be trusted. Three mechanisms independently drive the reported figure to zero without improving any
student's timetable:

\begin{enumerate}
  \item \textbf{Re-sectioning.} If student switching and initial sectioning are enabled, UniTime
    moves students between parallel laboratory batches until the overlaps disappear. The timetable
    is unchanged; the cohort structure of the institution is silently discarded.
  \item \textbf{An empty student population.} Under the \emph{Projected Student Course Demands}
    setting, and with no last-like enrolment data in the session, the solver loads zero students and
    reports a conflict count of zero over an empty denominator.
  \item \textbf{A denominator mismatch.} A conflict rate computed over the students who carry
    enrolments in a solver export is not comparable to one computed over every simulated student in
    the source data, and the two populations differed by several hundred students across the
    instances used in this project.
\end{enumerate}

The problem addressed in Sections~\ref{sec:evaluation} and~\ref{sec:experiments} is therefore stated
as two objectives rather than one: produce a feasible timetable, and establish a measurement of
student conflicts that cannot be satisfied by any of the three mechanisms above.
